\documentclass[11pt]{article}
\usepackage[margin=1in]{geometry}
\usepackage{amsmath,amssymb,amsthm,booktabs}
\newtheorem{theorem}{Theorem}
\newtheorem{lemma}{Lemma}
\newcommand{\E}{\mathbf E}
\newcommand{\Pp}{\mathbf P}
\newcommand{\one}{\mathbf 1}
\newcommand{\cF}{\mathsf F}
\newcommand{\cT}{\mathsf T}
\title{A Rooted-Grove Portmanteau Formula for Finite Markov-Chain Occupation Times}
\author{ProofNet24 P02 closure package for Aldous--Fill Open Problem 2.34}
\date{17 September 2026}

\begin{document}
\maketitle

\begin{abstract}
Aldous and Fill ask in Open Problem 2.34 whether the occupation-time identities of their Section 2.2.2 admit a single theorem and a Feynman-diagram-like combinatorial rule in the marked stationary/fundamental-matrix variables.  We give such a portmanteau formula for every finite irreducible discrete-time Markov chain, without reversibility.  The intrinsic combinatorial objects are directed spanning forests converging to a prescribed boundary.  A killed Green entry is a rooted-grove partition-function ratio.  Independently, the same ratio equals a bordered determinant in the Aldous--Fill data $(\pi,Z)$ with the denominator
\[
\Delta_A=\det\begin{pmatrix}Z_A^T&-\pi_A\\ \one^T&0\end{pmatrix}
\]
retained throughout.  An arbitrary stopping-time/boundary-data version contains all eleven identities in Section 2.2.2 as specializations.  The constituent occupation-measure and forest/Green mechanisms have substantial classical prior art; no publication-level novelty claim is made for those ingredients, and the exact-overlap question for the complete package remains a literature-classification issue.
\end{abstract}

\section{Source scope and notation}
Let $P=(p_{uv})_{u,v\in V}$ be the transition matrix of a finite irreducible discrete-time Markov chain.  Reversibility is not assumed.  Let $\pi$ be its stationary law, let
\[
L=I-P,\qquad \Pi=\one\pi^T,
\]
and let $Z$ be the Aldous--Fill fundamental matrix
\[
Z=\sum_{t\ge0}(P^t-\Pi),
\]
with the Ces\`aro (equivalently continuized) interpretation in the periodic case.  Equivalently,
\[
Z=(L+\Pi)^{-1}-\Pi,
\qquad LZ=ZL=I-\Pi,
\qquad Z\one=0,
\qquad \pi^TZ=0.
\]
For a stopping time $S$, ``visits before $S$'' includes time $0$ and excludes time $S$, matching the source convention.

For nonempty $A\subset V$, set $U=V\setminus A$ and $T_A=\inf\{t\ge0:X_t\in A\}$.  Write $Z_A$ for the principal submatrix indexed by $A$, and $\pi_A$ for the column vector $(\pi_a)_{a\in A}$.  The denominator frozen in the P02 program is
\begin{equation}
M_A=\begin{pmatrix}Z_A^T&-\pi_A\\ \one^T&0\end{pmatrix},
\qquad
\Delta_A=\det M_A.
\tag{1.1}
\end{equation}

\section{Rooted forests}
An $A$-rooted converging spanning forest is a spanning directed forest in which every $a\in A$ has outdegree $0$, every $u\notin A$ has outdegree $1$, and following arrows from every nonroot eventually reaches $A$.  Its weight is the product of the transition probabilities on its directed edges.  Let
\[
\cF_A=\sum_{F:\,A\text{-rooted}}\prod_{(u\to v)\in F}p_{uv}.
\]
Let
\[
\cT=\sum_{r\in V}\cF_{\{r\}}
\]
be the total weight of all one-root converging spanning trees.  For $i,j\notin A$, let
\[
\cF_{A\cup\{j\}}(i\to j)
\]
denote the total weight of forests rooted at $A\cup\{j\}$ in which $i$ belongs to the component rooted at $j$.

The directed all-minors matrix-tree theorem gives
\begin{equation}
\det L_{UU}=\cF_A,
\qquad
C_{ji}(L_{UU})=\cF_{A\cup\{j\}}(i\to j),
\tag{2.1}
\end{equation}
where $C_{ji}$ is the cofactor occurring in the $(i,j)$ entry of $L_{UU}^{-1}$.

\section{The portmanteau theorem}
For a probability law $\sigma$ on $V$, write
\[
z_{\sigma j}=\sum_{v\in V}\sigma_v Z_{vj},
\qquad
z_{\sigma,A}=(z_{\sigma a})_{a\in A}.
\]

\begin{theorem}[rooted-grove occupation portmanteau]
\label{thm:port}
Let $P$ be a finite irreducible discrete-time chain and $A\ne\varnothing$.

\emph{(a) First-hit occupation.} For $i,j\in U$,
\begin{equation}
\E_i\sum_{t<T_A}\mathbf 1_{\{X_t=j\}}
=
\frac{\cF_{A\cup\{j\}}(i\to j)}{\cF_A}.
\tag{3.1}
\end{equation}
Moreover
\begin{equation}
\boxed{\Delta_A=\frac{\cF_A}{\cT}>0.}
\tag{3.2}
\end{equation}
Define the frozen bordered numerator
\begin{equation}
\mathcal N_{A;i,j}
=
\det\begin{pmatrix}
Z_A^T&-\pi_A&Z_{iA}^T\\
\one^T&0&1\\
Z_{Aj}^T&-\pi_j&Z_{ij}
\end{pmatrix}.
\tag{3.3}
\end{equation}
Then
\begin{equation}
\boxed{\mathcal N_{A;i,j}=\frac{\cF_{A\cup\{j\}}(i\to j)}{\cT}},
\qquad
\boxed{\E_i\sum_{t<T_A}\mathbf 1_{\{X_t=j\}}=\frac{\mathcal N_{A;i,j}}{\Delta_A}}.
\tag{3.4}
\end{equation}

\emph{(b) Arbitrary stopping/boundary data.} Let $T$ be an integrable stopping time, let the initial law be $\sigma$, assume the terminal law $\tau$ is supported on $A$, and define the occupation row vector
\[
x_j=\E_\sigma\sum_{t<T}\mathbf 1_{\{X_t=j\}},
\qquad c=x_A.
\]
For $q\in U$ put
\begin{equation}
\mu_q=\sigma_q+\sum_{a\in A}c_a p_{aq}.
\tag{3.5}
\end{equation}
Then, for every $j\in U$,
\begin{equation}
\boxed{
x_j=\frac1{\cF_A}\sum_{q\in U}\mu_q\cF_{A\cup\{j\}}(q\to j).
}
\tag{3.6}
\end{equation}
Equivalently, if
\begin{equation}
\mathcal N_{\sigma,c;A,j}
=
\det\begin{pmatrix}
Z_A^T&-\pi_A&z_{\sigma,A}^T-c^T\\
\one^T&0&1\\
Z_{Aj}^T&-\pi_j&z_{\sigma j}
\end{pmatrix},
\tag{3.7}
\end{equation}
then
\begin{equation}
\boxed{x_j=\frac{\mathcal N_{\sigma,c;A,j}}{\Delta_A}},
\qquad
\boxed{\mathcal N_{\sigma,c;A,j}
=\frac1{\cT}\sum_{q\in U}\mu_q\cF_{A\cup\{j\}}(q\to j)}.
\tag{3.8}
\end{equation}
The boundary coordinates are, by definition, $x_a=c_a$.
\end{theorem}

\begin{proof}
Because the chain is finite and irreducible and $A$ is nonempty, the killed matrix $P_{UU}$ has spectral radius less than one.  Hence its Green matrix is
\[
G_A=(I-P_{UU})^{-1}=L_{UU}^{-1},
\]
and $(G_A)_{ij}$ is the expected number of visits to $j$ before $T_A$ when started from $i$.  Cramer's rule and (2.1) give
\[
(G_A)_{ij}
=\frac{C_{ji}(L_{UU})}{\det L_{UU}}
=\frac{\cF_{A\cup\{j\}}(i\to j)}{\cF_A},
\]
which proves (3.1).

We next identify the frozen denominator.  For $t\ne0$ put
\[
L_t=L+t\Pi.
\]
The identities for $Z$ imply
\begin{equation}
L_t^{-1}=Z+t^{-1}\Pi.
\tag{3.9}
\end{equation}
The Markov-chain tree theorem gives $\pi_r=\cF_{\{r\}}/\cT$, and consequently $\operatorname{adj}(L)=\cT\,\one\pi^T$.  The rank-one determinant formula therefore yields
\begin{equation}
\det L_t=t\cT.
\tag{3.10}
\end{equation}
Jacobi's complementary-minor identity applied to the invertible matrix $L_t$ gives
\begin{equation}
\det\bigl((L_t^{-1})_{AA}\bigr)
=
\frac{\det((L_t)_{UU})}{\det L_t}.
\tag{3.11}
\end{equation}
By (3.9), the left side equals
\[
\det\bigl(Z_A+t^{-1}\one_A\pi_A^T\bigr).
\]
The rank-one determinant identity and the transpose of (1.1) show that this is
\begin{equation}
\det Z_A+t^{-1}\Delta_A.
\tag{3.12}
\end{equation}
On the right of (3.11), (3.10) gives
\begin{equation}
\frac{\det(L_{UU}+t\one_U\pi_U^T)}{t\cT}.
\tag{3.13}
\end{equation}
Multiplying (3.11) by $t$ and letting $t\to0$ yields
\[
\Delta_A=\frac{\det L_{UU}}{\cT}=\frac{\cF_A}{\cT}>0,
\]
proving (3.2).

To obtain the marked numerator without defining the combinatorics from a determinant, use the occupation conservation identity.  For the first-hit problem, the terminal law is supported on $A$ and the boundary occupation is zero.  More generally, for part (b), pathwise flow conservation gives
\begin{equation}
x(I-P)=\sigma-\tau.
\tag{3.14}
\end{equation}
Restricting the columns of (3.14) to $U$ and using $\tau_U=0$ gives
\begin{equation}
x_U(I-P_{UU})=\sigma_U+cP_{AU}=\mu,
\tag{3.15}
\end{equation}
so
\[
x_j=\sum_{q\in U}\mu_q(G_A)_{qj}.
\]
Substitution of (3.1) proves the rooted-grove formula (3.6).

For the marked $(\pi,Z)$ form, let $s=\sum_v x_v=\E_\sigma T$.  Multiplying (3.14) by $Z$ and using $(I-P)Z=I-\Pi$ gives
\begin{equation}
x=\sigma Z-\tau Z+s\pi^T.
\tag{3.16}
\end{equation}
Restricting (3.16) to $A$ and using $\one^T\tau_A=1$ yields the bordered linear system
\begin{equation}
M_A\binom{\tau_A}{s}
=
\binom{z_{\sigma,A}^T-c^T}{1}.
\tag{3.17}
\end{equation}
Since (3.2) gives $\det M_A>0$, this system has a unique solution.  For $j\in U$, (3.16) says
\[
x_j=z_{\sigma j}-Z_{Aj}^T\tau_A+\pi_js.
\]
Taking the Schur complement of $M_A$ in the determinant (3.7) gives exactly
\[
\mathcal N_{\sigma,c;A,j}=\Delta_A x_j.
\]
This proves the first equality in (3.8).  Combining it with (3.2) and (3.6) gives the second.  Taking $\sigma=\delta_i$ and $c=0$ gives (3.3)--(3.4).  Thus the forest object is defined independently in the transition graph, and the bordered determinant is a second evaluation of the same partition-function ratio.
\end{proof}

\section{The source prototype}
For four distinct marked states $i,j,k,l$, take $A=\{k,l\}$.  The theorem gives the completely intrinsic rule
\begin{equation}
\boxed{
\E_i\bigl(\#\text{ visits to }j\text{ before }\min(T_k,T_l)\bigr)
=
\frac{\cF_{\{k,l,j\}}(i\to j)}{\cF_{\{k,l\}}}.
}
\tag{4.1}
\end{equation}
The combinatorial structure is: $k,l,j$ are roots, $i$ is constrained to the component rooted at $j$, and every unmarked state is summed over as an internal forest vertex.  By (3.2)--(3.4), (4.1) is also the ratio of the frozen marked $(\pi,Z)$ numerator and denominator.  This is the requested ``count the same quantity twice'' mechanism rather than a definition obtained by expanding a determinant after the fact.

\section{All eleven Section 2.2.2 identities}
We now account for the source list explicitly.  The basic singleton hitting specialization of (3.16) is worth recording.  If $i\ne l$, start from $i$, stop at $T_l$, and put $A=\{l\}$, $c_l=0$.  Then for every $j\ne l$,
\begin{equation}
\E_i\sum_{t<T_l}\mathbf1_{\{X_t=j\}}
=Z_{ij}-Z_{lj}+\pi_j\E_iT_l.
\tag{5.1}
\end{equation}
Also the boundary equation at $l$ is
\begin{equation}
0=Z_{il}-Z_{ll}+\pi_l\E_iT_l.
\tag{5.2}
\end{equation}

\begin{enumerate}
\item \textbf{Lemma 2.5.}  Start at $i$ and stop at the strict return $T_i^+$.  Here $A=\{i\}$, $\tau_i=1$, and the boundary occupation is $c_i=1$ because time zero is counted.  Equation (3.16) becomes $x=\pi^Ts$, since the two $Z$ terms cancel.  As $x_i=1$, $s=\E_iT_i^+=1/\pi_i$.

\item \textbf{Lemma 2.6.}  In the same strict-return specialization, $x_j=\pi_j s=\pi_j/\pi_i$.

\item \textbf{Lemma 2.7.}  For $j\ne i$, use (5.1) with start $j$, boundary $i$, and target coordinate $j$:
\[
\E_j\sum_{t<T_i}\mathbf1_{\{X_t=j\}}
=Z_{jj}-Z_{ij}+\pi_j\E_jT_i.
\]
By (5.2) with start $i$ and boundary $j$, $Z_{jj}-Z_{ij}=\pi_j\E_iT_j$.  Hence the displayed expectation is $\pi_j(\E_jT_i+\E_iT_j)$.

\item \textbf{Corollary 2.8.}  Starting from $i$, the number of visits to $i$ before $T_j$ is geometric with success probability $\Pp_i(T_j<T_i^+)$.  Its mean, by the preceding item with $i,j$ interchanged, is $\pi_i(\E_iT_j+\E_jT_i)$.  Taking the reciprocal gives the source formula.

\item \textbf{Lemma 2.9.}  Formula (5.1) and (5.2) give
\[
Z_{ij}=Z_{jj}-\pi_j\E_iT_j,
\qquad
Z_{lj}=Z_{jj}-\pi_j\E_lT_j.
\]
Substitution into (5.1) yields
\[
\E_i\sum_{t<T_l}\mathbf1_{\{X_t=j\}}
=\pi_j(\E_iT_l+\E_lT_j-\E_iT_j),
\]
including the source's endpoint conventions.

\item \textbf{Corollary 2.10.}  For $i\ne l$ and $j\ne l$, the strong Markov property gives
\[
\E_i\sum_{t<T_l}\mathbf1_{\{X_t=j\}}
=\Pp_i(T_j<T_l)\,
\E_j\sum_{t<T_l}\mathbf1_{\{X_t=j\}}.
\]
Apply Lemma 2.9 to the two expectations and divide to obtain exactly
\[
\Pp_i(T_j<T_l)
=\frac{\E_iT_l+\E_lT_j-\E_iT_j}{\E_jT_l+\E_lT_j}.
\]

\item \textbf{Lemma 2.11.}  Start with $\sigma=\pi$ and stop at $T_i$.  The boundary occupation is $c_i=0$.  Since $\pi^TZ=0$, the $i$-coordinate of (3.16) gives
\[
0=-Z_{ii}+\pi_i\E_\pi T_i,
\]
so $\pi_i\E_\pi T_i=Z_{ii}$.

\item \textbf{Lemma 2.12.}  Equation (5.2), with boundary $j$, is precisely
\[
\pi_j\E_iT_j=Z_{jj}-Z_{ij}.
\]

\item \textbf{Corollary 2.13.}  Sum Lemma 2.12 over $j$ and use $\sum_j Z_{ij}=0$:
\[
\sum_j\pi_j\E_iT_j=\sum_j Z_{jj}.
\]

\item \textbf{Corollary 2.14 (random target lemma).}  The right side in Corollary 2.13 is independent of $i$.

\item \textbf{Lemma 2.15.}  Again start from $\pi$ and stop at $T_i$.  For $j\ne i$, (3.16) and $\pi^TZ=0$ give
\[
\E_\pi\sum_{t<T_i}\mathbf1_{\{X_t=j\}}
=-Z_{ij}+\pi_j\E_\pi T_i
=\frac{\pi_j}{\pi_i}Z_{ii}-Z_{ij},
\]
where Lemma 2.11 was used in the last step.  The $j=i$ case is $0=0$ under the source convention.
\end{enumerate}
Thus one boundary-data theorem with the one frozen denominator $\Delta_A$ contains all eleven source identities.

\section{Prior art and classification}
The mathematical mechanisms used above are not presented as new in isolation.  Pitman's 1977 work on occupation measures gives a basic identity unifying many Markov-chain moment formulas.  Elementary finite-state matrix treatments predate the Aldous--Fill open problem.  The directed matrix-tree and all-minors machinery is classical.  Chebotarev and Agaev (2002) develop forest representations of Laplacian generalized inverses and eigenprojections.  Pitman and Tang (2018) give tree/forest formulas for finite-chain Green functions, hitting probabilities, mean first-passage times and Kemeny's constant.  Chung and Zeng (2023) give forest formulas for discrete Green functions on directed and weighted graphs.

Accordingly, the claims here are deliberately narrow.  The source-level answer is the single boundary-data compiler together with the exact bridge
\[
\text{rooted-grove partition function}
\longleftrightarrow
\text{frozen marked }(\pi,Z)\text{ bordered determinant},
\]
and the explicit accounting of the eleven Section 2.2.2 identities.  A bounded search through the sources above and direct searches for Open Problem 2.34 did not locate this exact complete package.  Failure to locate it is not a novelty certificate.  Publication-level novelty therefore remains unverified.

\section{Audit and certificate record}
The canonical theorem claim is `P02-OCC-C5-ROOTED-GROVE-PORTMANTEAU`, in the P02 report created at commit
\[
\texttt{8d40b5b3f5370e8209f63015a19b8a957be10e57}.
\]
Two independent end-to-end audits by other tracks are recorded as
\begin{itemize}
\item `P26-AUDIT-P02-I0023-C5-ROOTED-GROVE-04`, creation commit `12bb57d37208bfc01baf494e46f6a87a4f1c90e4`;
\item `P01-AUDIT-P02-I0023-C5-ROOTED-GROVE-74`, creation commit `77d2229b99e1eee42e0349facfbab14205514fac`.
\end{itemize}
P26 reconciles them under `P26-RECONCILE-P02-I0023-C5-TWO-AUDITS-07` as two distinct mathematical/source audits.  The accompanying `check.py` and `check.out` provide the archived exact-rational finite certificate for the forest/determinant bridge; the certificate includes reversible, periodic nonreversible and nonreversible chains.

\begin{thebibliography}{9}
\bibitem{AF}
D. Aldous and J. A. Fill,
\emph{Reversible Markov Chains and Random Walks on Graphs}, unfinished manuscript (2002), 2014 recompilation, Chapter 2, Open Problem 2.34.

\bibitem{Pitman1977}
J. W. Pitman,
Occupation measures for Markov chains,
\emph{Advances in Applied Probability} 9 (1977), 69--86.

\bibitem{CA2002}
P. Chebotarev and R. Agaev,
Forest matrices around the Laplacian matrix,
\emph{Linear Algebra and its Applications} 356 (2002), 253--274.

\bibitem{PT2018}
J. Pitman and W. Tang,
Tree formulas, mean first passage times and Kemeny's constant of a Markov chain,
\emph{Bernoulli} 24 (2018), 1942--1972.

\bibitem{CZ2023}
F. Chung and J. Zeng,
Forest formulas of discrete Green's functions,
\emph{Journal of Graph Theory} 102 (2023), 556--577.
\end{thebibliography}

\end{document}
